\documentclass[12pt,letterpaper,boxed]{hmcpset}
\usepackage[margin=1in,headheight=14pt]{geometry}
\usepackage{amsfonts, amsmath, amssymb, enumerate, fancyhdr, gensymb, lastpage, mathtools, parskip, graphicx}
\usepackage{xcolor, tikz-cd}
\newcommand{\wg}[1]{\textcolor{violet}{#1}}
\newcommand{\OO}{\mathcal O}
\newcommand{\Q}{\mathbb Q}
\newcommand{\R}{\mathbb R}
\newcommand{\C}{\mathbb C}
\newcommand{\Z}{\mathbb Z}
\renewcommand{\H}{\mathbb H}
\newcommand{\Hbdy}{\overline \H \cup \infty}
\newcommand{\abs}[1]{\left|#1\right|}
\newcommand{\im}{\text{im }}
\newcommand{\inv}{^{-1}}
\newcommand{\normal}{\unlhd} %% one can also use \trianglelelefteq
\newcommand{\anglee}[1]{\langle #1 \rangle}
\newcommand{\rest}[1]{\vert_{#1}}

\usepackage[shortlabels]{enumitem}

% Numbering macros
\pagestyle{fancy}
\lhead{Will Gilroy}
\chead{Algs Homework \#}
\rhead{03 November 2021}
\lfoot{}
\cfoot{}
\rfoot{Page\ \thepage\ of\ \pageref{LastPage}}

\linespread{1.5}

\newcommand\blankpage{
    \thispagestyle{empty}
    \addtocounter{page}{-1}
    \newpage}
\renewcommand\footrulewidth{0.4pt}

\begin{document}

\problemlist{Algorithms HW } 

%------------------------- Problem 1 -----------------------

\begin{problem}
	\includegraphics[scale=0.8]{1.png}
	\hfill
\end{problem}

\begin{solution}
Suppose for contradiction that a Riemann map $\Delta \to \mathcal C$
could not be extended continuously to the boundary. That is, we would
have two sequences in $\Delta$ converging to some point on the
boundary, $z_n, z_n' \to z_0 \in S^1$, such that their images
converged to distinct points on $\partial \mathcal C$. 
In other words, there is some positive $\delta$ such that $\abs{f(z_n)
- f(z_m')} > \delta$ for all $n, m$.

Consider the following series of calculation. 
Consider circles $\alpha_r$ of radius $r$ about $z_0$. For sufficiently large
$n,m$ and sufficiently small $r$ there are points in both sequences
which lie on $\alpha_r$. Now, by the fundamental theorem of calculus
we have 
\begin{align*}
	\delta < \abs{f(z_n) - f(z_m')} &= \abs{ \int_{\alpha_r} f'(\xi) d\xi} \\
		&= \abs{\int_{\theta_1(r)}^{\theta_2(r)} f'(z_0 + re^{i \theta}) r d\theta} \\
		&\leq \abs{\int_{\theta_1(r)}^{\theta_2(r)}} \abs{f'(z_0 + re^{i\theta})} r d\theta
\end{align*}
where the second equality is given by parameterizing $\alpha_r$ and
the final inequality given by the triangle inequality; the
angles $\theta_1(r), \theta_2(r)$ are the angles on which $z_n, z_m'$ respectively
lie on $\alpha_r$.
Now, squaring both sides and then applying Cauchy-Schwarz (taking 
$h = f r^{1/2}$, $g = r^{1/2}$) gives
\begin{align*}
	\delta^2 
	&\leq \int \abs{f'(z)}^2 r d\theta \int_{\theta_1(r)}^{\theta_2(r)} r d\theta \\
	&\leq \int \abs{f'(z)}^2 r d\theta \int_{0}^{2\pi} r d\theta  \\
	&\leq \int \abs{f'(z)}^2 r d\theta \left(2\pi r\right).
\end{align*}
Rearranging this inequality and then integrating over $r$ gives 
\[
	\int_{r=0}^{R} \frac{\delta^2}{r} dr 
	\leq \int_{r} \int_{\theta} \abs{f'(z)}^2 r d\theta dr = \int\int \abs{f'}^2  dA = \text{Area}(f(\Delta)),
\]
where the last equality is given by the fact that the Jacobian for the
appropriate transformation is $\abs{f'(z)}^2$. 
Now finally, notice that the left hand side diverges, but the right
hand side is finite, since it is the area of a bounded open set. A
contradiction, and so it must have been that $f$ extended continuously
to the boundary.

\wg{understand this jacobian statement, make a comment about the boundary being nice. }

\end{solution}

\newpage

%------------------------- Problem 2 -----------------------

\begin{problem}
	\includegraphics[scale=0.8]{2.png}
	\hfill
\end{problem}

\begin{solution}
Recall that we have a biholomorphism from the $\Delta \to \H$ given by 
\[
	h(z) = i \cdot \frac{1 + z}{1 - z}.
\]
This map extends continuously to the boudaries $\bar \Delta \to \Hbdy$
with $h(1) = \infty$. 
And so we can define a holomorphic function $g: \H \to \H$ by 
$g = h \circ f \circ h\inv$ which extends continuously to the actual
boundary $\bar \H \cup \infty$ and $g(\bar \H \cup \infty) \subseteq
\bar \H \cup \infty$.
The preimages of $1$ under $f$ are those which get sent to infinity
under $g$. 

Suppose for a moment that we have a single point on the real line
which gets sent to infinity. Say that we translate our coordinates so
that this point is $z=0$. Let us denote the ``right half upper plane''
by $\H^+ := \{z \in \H : Re(z) > 0\}$ and similarly let $\H^-$ denote
the ``left upper half plane''.

Now, the restriction of $g$ to $\H^+$ defines a holomorphic function
which extends continuously to the interval $(0, \infty)$ with the
image of $(0, \infty) \subseteq \R$. 
The Schwarz reflection principle then allows us to extend
$g\rest{\H^+}$ to a holomorphic function on the entire right half
plane $\R^+ := \{z \in \C: Re(z) > 0 \}$. The same argument means that
we can extend $g\rest{\H^-}$ to a holomorphic function on the entire left half plane.
We now have a holomorphic function $\C \setminus L \to \C$ where $L$
is the vertical line through $z=0$. Let us denote this function by
$\tilde g$. 

\begin{center}
	\begin{tikzpicture}[scale=2]

  % Axes
  \draw[->] (-2,0) -- (2,0) node[right] {$Re(z)$};
  %\draw[->] (0,-2) -- (0,2) node[above] {$\Im(z)$};

  % Dotted vertical line at z = 0
  \draw[dotted, thick] (0,-2) -- (0,2);

  % Labels for regions
  \node at (1,1) {$\mathbb{H}^+$};
  \node at (-1,1) {$\mathbb{H}^-$};
  \node at (0.3,-0.2) {$z=0$};
  \fill (0,0) circle (1pt);
\end{tikzpicture}
\end{center}

We can extend $\tilde g$ further to a function $G: \C \setminus
\{z=0\} \to \C$ in the following way
\[
	G(z) = \begin{cases}
		\tilde g(z) & z \not\in L \\
		g(z) & z \in L \cap \H \\
		\overline{g(\bar z)} & z \in L \cap -\H.
	\end{cases}
\]
The proof that Scharz reflection gives a holomorphic function also
proves that this function is holomorphic as a function $\C \setminus
\{z = 0 \} \to \C$. 

We can translate $G$ into a holomorphic extension of $f$ by
conjugating $h\inv \circ G \circ h : \C \setminus \{p\} \to \C$ where
$p$ is the preimage of $1$ under the continuous extension $\overline
f$.

If $f$ has more than one point on $S^1$ which is sent to $1$ (i.e.
more than one point on $\R$ is sent to $\infty$ under $h$) then we can
repeat a similar procedure, except now we Schwarz reflect vertical
strips bounded by the vertical lines through the points on the real
line which are sent to $\infty$ by $h$.

\wg{Lastly, must there only be finitely many points?}









\end{solution}

\newpage

%------------------------- Problem 3 -----------------------

\begin{problem}[3]
	\hfill
\end{problem}
\begin{solution}
\end{solution}

\newpage

%------------------------- Problem 4 -----------------------

\begin{problem}
	\includegraphics[scale=0.8]{4.png}
	\hfill
\end{problem}

\begin{solution}
We know from the Riemann mapping theorem that any simply connected open set 
(which is not all of $\C$) is biholomorphic to the unit disc. And so,
any two bounded simply connected open sets are biholomorphic to each other.
Thus, we need at least one of our open sets to be non-simply connected.
Moreover, since we want to find a homeoomrphism between our open sets, it
follows then that both are not simply connected (and much have the same
homology). 
(From here, how do we show that there's no biholomorphism between two given open
sets)
\end{solution}

\newpage

%------------------------- Problem 5 -----------------------

\begin{problem}[4]
	\hfill
\end{problem}

\begin{solution}
\end{solution}

\newpage


\end{document}
